% ===========================================================================
%  第15章：级联姿态-角速度控制架构
% ===========================================================================

\section{级联姿态-角速度控制架构：解耦设计与交叉项补偿}
\label{sec:cascade_chapter}

前文分别发展了三种姿态控制架构：第\ref{sec:control}节的级联SAS（对角分配）、
第\ref{sec:INDI}节的INDI（在线效能矩阵+角加速度反馈）、
第\ref{sec:lqr}节的LQR（工作点最优全状态反馈）。
第\ref{sec:coupling_rejection}节与第\ref{sec:lqr_coupling_limit}节的耦合抑制对比表明，
三者在处理本构型固有的反扭矩投影交叉耦合时各有局限：
SAS的对角分配结构性放弃交叉项，INDI受角加速度估计噪声与增量饱和制约，
LQR的冻结效能矩阵随工作点漂移而退化。

本章提出并系统阐述\textbf{级联姿态-角速度控制架构}——这也是实机飞控（2026-08 状态）
最终采用的架构。与前三者的根本区别在于\textbf{时间尺度分离}：将控制问题分解为
带宽分离的两个回路，姿态环（慢）只负责姿态误差调节、角速度环（快）承担
交叉项补偿，使耦合扰动优先在内环受到抑制，减少其向姿态误差的积分传播。

\textbf{实机架构总览（五层控制链路，\texttt{flight\_control.cpp}）}：
\begin{equation}
\begin{aligned}
    \underbrace{\mathbf q_{\mathrm{target}}}_{\text{参考姿态}}
    &\xrightarrow{\;\mathbf q_{\mathrm{err}}=\mathbf q^{-1}\otimes\mathbf q_{\mathrm{target}}\;}
    \underbrace{\boldsymbol\omega_{\mathrm{ref}}}_{\text{四元数外环}}
    \xrightarrow{\;\text{速率PI}\;}
    \underbrace{\boldsymbol\alpha_{\mathrm{cmd}}^{\deg}}_{\text{角加速度指令（deg/s$^2$）}},\\
    \boldsymbol\alpha_{\mathrm{cmd}}^{\deg}
    &\xrightarrow{\;\boldsymbol\alpha_{\mathrm{cmd}}=(\pi/180)\boldsymbol\alpha_{\mathrm{cmd}}^{\deg}\;}
    \xrightarrow{\;\mathbf M=\mathbf I\boldsymbol\alpha_{\mathrm{cmd}}+\mathbf M_{\mathrm{ff}}\;}
    \underbrace{\mathbf M_{\mathrm{cmd}}}_{\text{惯量逆解与前馈}}
    \xrightarrow{\;\mathbf B_{\mathrm{true}}^{-1}(\mathbf u_{k-1})\;}
    \underbrace{[\Delta\omega,\delta_t,\delta_f]^T}_{\text{在线分配}} .
\end{aligned}
    \label{eq:cascade_chain}
\end{equation}
五层各自职能单一、参数独立，这正是"用几何换算法简洁、用结构换性能"设计哲学的
控制侧体现。本章按此链路组织：§\ref{sec:cascade_design}给出控制律设计
（含惯量逆解前馈与工作点管理），§\ref{sec:cascade_tf}--\ref{sec:cascade_sp}
给出线性化/灵敏度/奇异摄动理论，§\ref{sec:cascade_tuning}给出离线整定，
§\ref{sec:cascade_sim}给出仿真验证，最后§\ref{sec:cascade_implementation}
给出实机实现状态与工程教训。

\subsection{设计思想：时间尺度分离与功能解耦}
\label{sec:cascade_philosophy}

单级控制架构的根本困难在于：姿态误差动力学、角速度动力学与执行器交叉耦合
被压缩在同一个反馈回路中，控制律必须同时兼顾跟踪品质、稳定性与耦合补偿，
三者相互制约。第\ref{sec:lqr_coupling_limit}节保留的历史原型表明，效能矩阵在线化
（增益调度、SDRE或在线分配）并未在该组模型与权重下自动消除度级残差；
该现象提示需要分离姿态调节与快速耦合补偿，但不能据此断言所有单级全状态反馈均存在同一限制。

级联架构的设计思想是\textbf{用时间尺度分离换取功能解耦}：
将控制问题分解为两个带宽分离的回路，每个回路只承担单一职能。

\begin{itemize}
    \item \textbf{外环（姿态环，慢环）}：仅负责姿态误差的调节，输出角速度指令。
          由于不直接操作执行器，外环无需考虑效能矩阵、执行器饱和与交叉耦合，
          可采用极简的纯比例（P）或比例-积分（PI）结构；
    \item \textbf{内环（角速度环，快环）}：仅负责跟踪角速度指令，输出执行器指令。
          由于角速度环带宽显著高于姿态环（时间尺度分离假设），
          内环可在姿态还未来得及响应前快速抑制交叉耦合扰动，
          使其无法经角速度-姿态积分传播到姿态环。
\end{itemize}

这一思想的物理本质是：\textbf{交叉耦合的抑制速度必须快于它影响姿态的速度}。
在单级架构中，补偿力矩与耦合扰动同时作用于同一动力学，无法形成时间尺度上的先后；
级联架构则通过内环的高带宽，使耦合扰动在姿态环"察觉"之前即被消除。
这与第\ref{sec:control}节SAS的级联思想一脉相承，但内环用在线效能矩阵
替代对角近似，形成工作点相关的准线性分配；它仍不是非线性映射的全局精确逆。

\subsection{控制律设计}
\label{sec:cascade_design}

\subsubsection{外环：误差四元数几何控制与精确角度提取}

姿态反馈采用第\ref{sec:quat_control}节的误差四元数，
以规避欧拉角在$\theta=\pm90^\circ$处的万向锁奇异性。
设当前姿态四元数$\mathbf{q}$、期望姿态$\mathbf{q}_{\mathrm{target}}$，
误差四元数$\mathbf{q}_{\mathrm{err}}=\mathbf{q}^{*}\otimes\mathbf{q}_{\mathrm{target}}$
的矢量部分$\boldsymbol{\varepsilon}_{\mathrm{err}}\in\mathbb{R}^3$编码姿态误差。
实机（\texttt{flight\_control.cpp} \texttt{execute\_attitude\_controller}）的
外环实现分三步：

\begin{enumerate}[label=(\arabic*)]
    \item \textbf{短弧归一化}：当$q_w<0$时将整个误差四元数乘以$-1$，使$\eta_{\mathrm{err}}\geq0$
          ——四元数双覆盖（$\mathbf{q}$ 与 $-\mathbf{q}$ 同姿态）下确保沿最短路径旋转；
    \item \textbf{精确角度提取}（非小角度近似）：
          \begin{equation}
              \varepsilon_i = \mathrm{sign}(q_w)\cdot q_{v,i}\cdot
              \frac{2\arctan(\|\mathbf{q}_v\|/|q_w|)}{\|\mathbf{q}_v\|}
              \label{eq:cascade_angle_extract}
          \end{equation}
          多轴同时存在误差时 $\|\mathbf{q}_v\|$ 取\textbf{全模长}（非单轴投影），
          避免偏航误差分量使滚转/俯仰的 atan2 判断偏低；
    \item \textbf{通道映射}（悬停构型语义）：$\varepsilon_x\to$上摆（滚转）、
          $\varepsilon_y\to$下摆（俯仰）、$\varepsilon_z\to$航向（差速，见§\ref{sec:cascade_yaw}）。
\end{enumerate}

角速度指令由姿态误差经比例控制生成（实机为纯 P，$K_i=K_d=0$）：
\begin{equation}
    \boxed{
    \boldsymbol{\omega}_{\mathrm{ref}}
    = \mathbf{K}_p\,\boldsymbol{\varepsilon}_{\mathrm{err}}
    }
    \label{eq:cascade_outer_ch}
\end{equation}
其中滚转/俯仰外环$\mathbf{K}_p$的实机比例增益均为2.8（deg域）。航向不直接使用
$\varepsilon_z$的同构外环，而采用第\ref{sec:cascade_yaw}节的ratchet hold：杆量给定
角速度，回中后用短弧航向误差和增益5.0生成限幅角速度参考。
式(\ref{eq:cascade_angle_extract})输出的是完整旋转向量，量级近似为姿态角误差，
不是原始四元数矢量部分的半角。因此$K_p$直接对应姿态误差到角速度参考的局部斜率；
后续线性分析不能再额外引入$1/2$因子。

\subsubsection{内环：角速度环输出角加速度指令}

内环在角度域将角速度误差映射为\textbf{角加速度指令}：
\begin{equation}
    \boldsymbol{\alpha}_{\mathrm{cmd}}^{\deg}
    = \mathbf{K}_\omega^{\deg}\left(\boldsymbol{\omega}_{\mathrm{ref}}^{\deg}
                        - \boldsymbol{\omega}^{\deg}\right)+\boldsymbol\alpha_I^{\deg}
    \label{eq:cascade_nu_ch}
\end{equation}
实机为速率PI。2026-08-11量纲迁移后，滚转/俯仰
$K_p=16.042818$、航向$K_p=11.459156$，三轴$K_i=5.729578$；其输出、
输出限幅和积分限幅均以deg/s$^2$计。进入物理力矩层前，
\texttt{ControlUnits::dps2ToRadps2}只执行一次$\pi/180$转换，得到
$\boldsymbol\nu=\boldsymbol\alpha_{\mathrm{cmd}}$（rad/s$^2$）。这些数值由迁移前
0.28/0.20/0.10整体乘$180/\pi$得到，因此物理指令和执行器输出严格保持。
例如滚转/俯仰速率误差每增加$1\,\mathrm{deg/s}$，比例项增加
$16.042818\,\mathrm{deg/s^2}$；这就是当前$K_p$的直接物理含义，
不再需要额外解释为“乘57.3后的等效增益”。积分分离阈值仍是角速度误差
$60\,\mathrm{deg/s}$，不参与缩放。
积分项对未建模的慢扰动（气动交叉/重心偏置）提供稳态消除——
这是"反馈兜底"哲学的体现：显式对消（§\ref{sec:cascade_ff}）处理已知项，
积分项处理未知项。

\subsubsection{惯量逆解与交叉耦合前馈（NDI 式 $f(\mathbf{x})$ 对消）}
\label{sec:cascade_ff}

角加速度指令经惯量矩阵成为期望力矩，并叠加\textbf{显式前馈对消项}：
\begin{equation}
    \boxed{
    \mathbf{M}_{\mathrm{cmd}} = \mathbf{I}\boldsymbol{\nu}
    + \underbrace{\boldsymbol{\omega}\times(\mathbf{I}\boldsymbol{\omega})
      + \boldsymbol{\omega}\times\mathbf{h}_{\mathrm{rotor}}}_{\text{前馈 }
      \mathbf{M}_{\mathrm{ff}}}
    }
    \label{eq:cascade_ff}
\end{equation}
其中$\mathbf{h}_{\mathrm{rotor}}=J_p(w_f\hat{n}_f + w_t\hat{n}_t)$为前后转子的净角动量
（反向旋转部分抵消，悬停时沿推力轴）。
前馈项即 NDI 公式 $\mathbf{u}=\mathbf{B}^{-1}(\boldsymbol{\nu}-f(\mathbf{x}))$
中的 $f(\mathbf{x})$ 项（陀螺耦合与转子陀螺）——\textbf{显式计算并补偿已知非线性}，
而非让反馈环"扛"。
实机（\texttt{InertiaDecoupling.h}，2026-08-10 接入）以掩码控制（默认全开，
可在线 A/B 关闭）；前馈在机体系按物理轴计算、经轴置换进入模型系后与
$\mathbf{I}\boldsymbol{\nu}$同坐标系相加。显式关系为
\begin{equation}
\mathbf M_{\mathrm{cmd}}^P
=\mathbf R_{P\leftarrow F}\mathbf M_{\mathrm{cmd}}^F .
\label{eq:cascade_axis_permutation}
\end{equation}
省略该矩阵会把任务滚转/偏航与模型$M_{z'}/M_{x'}$混写。
气动与重力项未显式对消（模型不确定度大），由内环积分项兜底。

\subsubsection{NDI 理论基础：模型逆与在线效能矩阵}

第\ref{sec:cascade_ff}节的前馈对消源于经典\textbf{非线性动态逆（NDI）}。考虑
输入仿射非线性系统 $\dot{\boldsymbol{\omega}} = \boldsymbol{\alpha}(\mathbf{x}) +
\mathcal{B}(\mathbf{x})\mathbf{u}$（$\boldsymbol{\alpha}$ 含 Coriolis 力、陀螺力矩、
气动阻尼等与 $\mathbf{u}$ 无关的动力学项），NDI 直接构造逆映射\cite{stevens2015aircraft}：
\begin{equation}
    \mathbf{u} = \hat{\mathcal{B}}^{-1}(\mathbf{x})\left[\boldsymbol{\nu}
        - \hat{\boldsymbol{\alpha}}(\mathbf{x})\right]
    \label{eq:ndi_law}
\end{equation}
其中 $\hat{\boldsymbol{\alpha}}$、$\hat{\mathcal{B}}$ 为机载模型估计，$\boldsymbol{\nu}$
为虚拟控制量（期望角加速度）。式(\ref{eq:cascade_ff})与NDI模型补偿思想相关，但当前绝对
Jacobian分配并非非线性映射的精确逆：
$\boldsymbol{\nu}$ 来自速率内环、$\mathbf{I}\boldsymbol{\nu}$ 为惯量逆、
前馈 $\mathbf{M}_{\mathrm{ff}}$ 对消 $\boldsymbol{\omega}\times(\mathbf{I}\boldsymbol{\omega})
+\boldsymbol{\omega}\times\mathbf{h}$（$\boldsymbol{\alpha}$ 中可精确建模的部分），
$\mathbf B_{\mathrm{true}}^{-1}$只提供上一拍局部增益。只有模型、工作点和坐标一致且高阶残差
可忽略时，闭环才近似呈$\dot{\boldsymbol\omega}\approx\boldsymbol\nu$。

\subsubsection{在线局部效能矩阵（B\_true）}
\label{sec:cascade_Btrue}

NDI 控制律中唯一需要建模的是控制效能矩阵
$\mathbf{B}=\partial\mathbf{M}/\partial\mathbf{u}$。对本构型，
控制输入$\mathbf u=[\Delta\omega,\delta_t,\delta_f]^T$，模型系力矩
$\mathbf M^P=[M_{x'},M_{y'},M_{z'}]^T$，从第\ref{sec:mapping}节的完整六维映射
（式\ref{eq:six_component}）的力矩分量出发：
\begin{equation}
    \begin{aligned}
        M_{x'} &= -\tau_f \cos\delta_f + \tau_t \cos\delta_t \\
        M_{y'} &= -b T_t \sin\delta_t - \tau_f \sin\delta_f \\
        M_{z'} &= a T_f \sin\delta_f - \tau_t \sin\delta_t
    \end{aligned}
    \label{eq:moment_full_nonlinear}
\end{equation}
其中推力与反扭矩由差速变量 $\Delta\omega$ 参数化（式\ref{eq:diff_def}）：
\begin{equation}
    T_f = T_0(1+\Delta\omega), \quad T_t = T_0(1-\Delta\omega), \quad
    \tau_f = \tau_0(1+\Delta\omega), \quad \tau_t = \tau_0(1-\Delta\omega)
    \label{eq:diff_thrust_torque}
\end{equation}
且 $T_0=k_T\omega_0^2$、$\tau_0=k_Q\omega_0^2$。

对式(\ref{eq:moment_full_nonlinear})逐项求偏导，得到\textbf{在线 Jacobian}
（在当前工作点 $\mathbf{u}_k=[\Delta\omega_k,\delta_{t,k},\delta_{f,k}]^T$ 处）：
\begin{equation}
    \boxed{
    \mathbf{B}_{\mathrm{true}}(\mathbf{u}_k) =
    \begin{bmatrix}
        -\tau_0(\cos\delta_{f,k}+\cos\delta_{t,k}) & -\tau_{t,k}\sin\delta_{t,k} & +\tau_{f,k}\sin\delta_{f,k} \\[6pt]
        +b T_0\sin\delta_{t,k} - \tau_0\sin\delta_{f,k} & -b T_{t,k}\cos\delta_{t,k} & -\tau_{f,k}\cos\delta_{f,k} \\[6pt]
        +a T_0\sin\delta_{f,k} + \tau_0\sin\delta_{t,k} & -\tau_{t,k}\cos\delta_{t,k} & +a T_{f,k}\cos\delta_{f,k}
    \end{bmatrix}
    }
    \label{eq:B_true}
\end{equation}
其中 $T_{f,k}=T_0(1+\Delta\omega_k)$ 等为当前步推力和反扭矩。

当 $\delta_f=\delta_t=0$ 时式(\ref{eq:B_true})退化为第\ref{sec:allocation}节的
线性化矩阵（式\ref{eq:effectiveness_matrix}）：
\begin{equation}
    \mathbf{B}_{\mathrm{true}}(\mathbf{0}) =
    \begin{bmatrix}
        -2\tau_0 & 0 & 0 \\
        0 & -b T_0 & -\tau_0 \\
        0 & -\tau_0 & a T_0
    \end{bmatrix}
    \equiv \mathbf{B}_{\mathrm{full}}(\omega_0)
    \label{eq:B_at_zero}
\end{equation}
即 $\mathbf{B}_{\mathrm{full}}$ 是 $\mathbf{B}_{\mathrm{true}}$ 的零摆角特例。
大摆角工况下两者差异显著（表\ref{tab:B_comparison}）——俯仰通道效能估计误差
超过 50\%，若用冻结 $\mathbf{B}_{\mathrm{full}}$ 求逆，俯仰指令会被放大近两倍：
\textbf{这正是实机必须采用在线 $\mathbf{B}_{\mathrm{true}}$ 的理论依据}。

\begin{table}[H]
\centering
\caption{$\mathbf{B}_{\mathrm{full}}$与$\mathbf{B}_{\mathrm{true}}$在MAXDEF工况下的对比
（$\delta_f=\delta_t=15^\circ$、$\Delta\omega=0.5$）}
\label{tab:B_comparison}
\setlength{\tabcolsep}{4pt}
\small
\begin{tabularx}{\textwidth}{c c c X}
\toprule
\textbf{矩阵元素} & $\mathbf{B}_{\mathrm{full}}(\omega_0)$ & $\mathbf{B}_{\mathrm{true}}(\mathbf{u}_k)$ & \textbf{相对误差} \\
\midrule
$B_{11}$（$M_{x'}\leftarrow\Delta\omega$） & $-2\tau_0$ & $-1.932\tau_0$ & 3.4\% \\
$B_{22}$（$M_{y'}\leftarrow\delta_t$） & $-bT_0$ & $-0.483bT_0$ & \textbf{51.7\%} \\
$B_{33}$（$M_{z'}\leftarrow\delta_f$） & $+aT_0$ & $+1.449aT_0$ & \textbf{44.9\%} \\
$B_{12}$（$M_{x'}\leftarrow\delta_t$） & 0 & $-0.129\tau_0$ & 新出现 \\
$B_{13}$（$M_{x'}\leftarrow\delta_f$） & 0 & $+0.388\tau_0$ & 新出现 \\
\bottomrule
\end{tabularx}
\end{table}

$\mathbf{B}_{\mathrm{true}}$ 的逆经伴随矩阵法解析求取（行列式
\begin{equation}
    \det(\mathbf{B}) = b_{11}(b_{22}b_{33}-b_{23}b_{32}) - b_{12}(b_{21}b_{33}-b_{23}b_{31}) + b_{13}(b_{21}b_{32}-b_{22}b_{31})
    \label{eq:B_det}
\end{equation}
元素$b_{ij}$为当前步物理量的组合。零摆角名义点可解析证明满秩；
有限摆角下应在限定工作域检查条件数，不能仅凭局部子块宣称$\det\mathbf B>0$全域严格成立。
$\omega_0\to0$时矩阵趋于奇异，由第\ref{sec:cascade_workpoint}节的工作点下限与门控处理。

\subsubsection{在线效能矩阵分配（绝对准线性形式，非INDI）}
\label{sec:cascade_alloc}

力矩指令经\textbf{在线 Jacobian 逆}分配为执行器指令：
\begin{equation}
    \mathbf{u}_k = \mathbf{B}_{\mathrm{true}}^{-1}(\mathbf{u}_{k-1})\,
                   \mathbf{M}_{\mathrm{cmd}},
    \qquad \mathbf{u} = [\Delta\omega,\;\delta_t,\;\delta_f]^T
    \label{eq:cascade_alloc_firmware}
\end{equation}
$\mathbf{B}_{\mathrm{true}}(\mathbf{u}_{k-1})$ 为在\textbf{上一拍执行器工作点}
处求值的在线 Jacobian（式\ref{eq:B_true}）——非零摆角名义点，
使分配矩阵跟踪工作点变化（gain-scheduling 思想）。
概念仿真中曾考察\textbf{增量变体}（$\mathbf{u}_{k+1}=\mathbf{u}_k+\Delta\mathbf{u}$，
INDI 风格）：
\begin{equation}
    \mathbf{u}_{k+1} = \mathbf{u}_k
        + \mathbf{B}_{\mathrm{true}}^{-1}(\mathbf{u}_k)\,
          \mathbf{I}\,\boldsymbol{\nu}\,\Delta t
    \label{eq:cascade_alloc_ch}
\end{equation}
实机未采用增量形式（理由见下"架构定性"）——本式仅用于与 NDI 绝对形式的
理论对比（第\ref{sec:cascade_tf}节连续化分析）与 INDI 方案（第\ref{sec:INDI}章）的
联系。

\textbf{架构定性}：当前固件是基于上一拍Jacobian的\textbf{绝对准线性分配}，不是INDI。
式(\ref{eq:cascade_alloc_firmware})也不是非线性映射的Newton逆；严格局部Newton更新应包含
$\mathbf u_{k-1}+\mathbf B^{-1}[\mathbf M_{\mathrm{cmd}}-\mathbf M(\mathbf u_{k-1})]$残差项。
当前实现可称NDI型级联，但不应称为精确非线性动态逆。判别依据为：
\begin{enumerate}[label=(\arabic*)]
    \item \textbf{绝对指令形式}：$\mathbf{u}_k$ 由力矩直接分配得到，
          无增量累加 $\mathbf{u}_{k+1}=\mathbf{u}_k+\Delta\mathbf{u}$；
    \item \textbf{无实测角加速度}：$\boldsymbol{\nu}$ 由角速度误差经比例/积分生成，
          $\mathbf{B}_{\mathrm{true}}^{-1}$ 作静态分配而非动态逆，
          不使用 IMU 差分 $\dot{\boldsymbol{\omega}}_{\mathrm{IMU}}$。
\end{enumerate}
INDI 对模型误差的部分鲁棒性在实机中由\textbf{$\tau_m$模型观测器与反馈积分}补充，
但两者不等价于基于实测角加速度的增量动态逆
（见§\ref{sec:cascade_workpoint}）：模型预测器减小指令与估计工作点的相位失配，
积分项吸收部分慢变误差；二者不等价于基于实测角加速度的INDI，也不能宣称取得同等效果。

\subsubsection{工作点管理：$\tau_m$模型观测器、增益调度与良态化}
\label{sec:cascade_workpoint}

实机分配层有三项工作点管理（均为 2026-08 抖油门震荡三层定位的产物）：
\begin{enumerate}[label=(\arabic*)]
    \item \textbf{$\tau_m$一阶模型观测器}（\texttt{flight\_control.cpp} 层1 之前）：
          \begin{equation}
              w_{\mathrm{est}} \mathrel{+}=
              (w_{\mathrm{cmd}} - w_{\mathrm{est}})\cdot\frac{\Delta t}{\tau_m}
              \label{eq:cascade_obs}
          \end{equation}
          该状态由指令和$\tau_m=0.28$\,s开环预测，没有转速传感器校正，故应称模型估计而非
          “观测实际转速”。它可减小油门瞬态的相位失配，但不能保证有效增益恒定；
          当前实现中差速列仍使用输入结构的指令基准转速，摆角列使用经下限处理的估计转速；
    \item \textbf{差速增益调度}：低油门时，$1/w_0^2$归一化会放大差速指令。
          固件采用
          \begin{equation}
              g_s=\min\!\left\{\left(\frac{w_{0,\mathrm{eff}}}
              {w_{\mathrm{hover}}}\right)^2,1\right\}
          \end{equation}
          限制该效应，详见第\ref{sec:allocation}节；
    \item \textbf{良态工作点下限与零油门门控}：B 各项 $\propto w_0^2$、det $\propto w_0^6$，
          低油门为数值病态区（5\% 油门 det 比 50\% 小 6 个数量级）——
          工作点按 $w_{0,\mathrm{floor}}=0.6\,w_{\mathrm{hover}}$ 冻结；
          零油门（$<$5\%）强制执行器归中（防 B 良态化后舵机对姿态误差误响应）。
\end{enumerate}

\subsubsection{航向通道：ratchet hold（摇杆恒角速度 + 回中锁航向）}
\label{sec:cascade_yaw}

悬停/巡航的航向通道\textbf{不做姿态回中}——航向摇杆恒为角速度指令
（与 RATE\_MODE 语义一致，见第\ref{sec:rc_curve}节摇杆曲线），回中时以
\textbf{航向短弧误差}锁住当时航向：
\begin{equation}
    \omega_{\mathrm{ref},z} =
    \begin{cases}
        \mathrm{map}(\text{航向摇杆}, \pm512\text{us}, \pm\omega_{\max,z}), & |\text{杆量}| > 40\text{us} \\
        \mathrm{clip}\!\left(k_p^{\mathrm{yaw}}\cdot\Delta\psi_{\mathrm{short}}, \pm 60^\circ/\mathrm{s}\right), & \text{回中}
    \end{cases}
    \label{eq:cascade_yaw_hold}
\end{equation}
其中 $\Delta\psi_{\mathrm{short}}\in(-\pi,\pi]$ 为参考航向的短弧误差（回绕归一化），
实机 $k_p^{\mathrm{yaw}}=5.0$；解锁瞬间参考重置为当前航向防起飞猛转。
该设计绕开四元数航向回中在悬停万向锁附近的欧拉表示退化，且满足
"摇杆恒角速度 + 松杆即停（锁航向）"的操纵习惯（2026-08-09 实机 ratchet hold）。

\subsubsection{人机接口：FPV 摇杆曲线}

速率模式与航向通道的杆量到角速率映射采用Betaflight三参数模型。
三个参数分别为\texttt{rc\_rate}、\texttt{rc\_expo}和\texttt{rc\_super}，
其公式与默认值见第\ref{sec:rc_curve}节。
外环与分配层完全不感知曲线形状，曲线只改变 $\boldsymbol{\omega}_{\mathrm{ref}}$ 的来源
（摇杆曲线或姿态误差经$K_p$），从而将人机接口与控制律解耦。

\begin{remark}[固件实现：NDI 式级联（绝对指令 + 在线 B + 工作点观测）]
\label{rem:firmware_alloc}
实机控制链即式(\ref{eq:cascade_chain})五层：
\begin{equation}
    \boldsymbol{\alpha}_{\mathrm{cmd}}^{\deg}
    \xrightarrow{\;\times\,\pi/180\;}
    \boldsymbol{\alpha}_{\mathrm{cmd}}
    \xrightarrow{\;\times\,\mathbf{I}\;+\;\mathbf{M}_{\mathrm{ff}}\;}
    \mathbf{M}_{\mathrm{cmd}} \xrightarrow{\;\mathbf{B}_{\mathrm{true}}^{-1}(\mathbf{u}_{k-1})\;}
    \mathbf{u}_k = [\Delta\omega,\;\delta_t,\;\delta_f]^T
    \label{eq:cascade_alloc_firmware2}
\end{equation}
其中$\boldsymbol{\alpha}_{\mathrm{cmd}}^{\deg}$为速率内环（PI）直接输出的deg/s$^2$指令，
$\boldsymbol{\alpha}_{\mathrm{cmd}}$为经唯一单位边界换算后的rad/s$^2$物理量，
$\mathbf{I}=\operatorname{diag}(I_x,I_y,I_z)$为惯量矩阵，
$\mathbf{B}_{\mathrm{true}}(\mathbf{u}_{k-1})$为在\textbf{上一拍执行器工作点}
处求值的在线Jacobian。

\textbf{工程理由（为什么实机选 NDI 式而非 INDI）}：
\begin{enumerate}[label=(\arabic*)]
    \item 速率PI直接输出$\boldsymbol{\alpha}$，经$\mathbf{I}$缩放即为力矩量纲，
          无需维护执行器状态$\mathbf{u}_k$的增量累积历史；
    \item 规避INDI对$\dot{\boldsymbol{\omega}}_{\mathrm{IMU}}$数值微分的需求，
          消除高频噪声经$\mathbf{B}^{-1}$放大的路径（第\ref{sec:estimator_study}节）；
    \item 未建模非线性（气动交叉等）由内环PI的积分项作为"慢扰动"补偿，
          而非显式对消——概念阶段参数不确定度大，显式对消反而可能引入更大误差；
    \item 在线辨识模块（\texttt{TandemVec\_OnlineID.h}，遗忘因子RLS）
          监测惯量比$b$与总扰动$d$（实测 $b\approx0.35/0.36/0.90$、$d_{\mathrm{yaw}}\approx-7.36$，
          见第\ref{sec:online_id}章），自适应增益建议$K_p/\sqrt{b}$为\textbf{纯观测模式}——
          仅写入遥测量供地面站显示与离线复算，未写回PID增益。
\end{enumerate}

\textbf{实机配置（2026-08 调参态）}：姿态外环 PositionPID 纯P（$K_p=2.8/2.8$，
四元数精确角度提取输入）；速率内环PI在deg域使用
$K_p=16.042818/16.042818/11.459156$、三轴$K_i=5.729578$，进入物理层前
统一转换为rad/s$^2$；
惯量前馈掩码全开（可 A/B）；工作点 $w_{0,\mathrm{floor}}=0.6\,w_{\mathrm{hover}}$
+ 增益调度封顶 1.0 + 零油门门控；航向 ratchet hold（$k_p^{\mathrm{yaw}}=5.0$）。
RATE\_MODE 满杆角速率 333°/s（限幅 600，2026-08-11 由 80 解耦——见第\ref{sec:rc_curve}节）。
\end{remark}

\subsection{闭环动力学与传递函数分析}
\label{sec:cascade_tf}

为理解级联结构，必须从固件的绝对分配而不是历史增量式$\dot{\mathbf u}$出发。
在坐标置换已正确处理、局部矩阵与植物一致、无饱和且忽略执行器动态时，固件近似满足
\begin{equation}
\mathbf u=\mathbf B^{-1}\mathbf I\boldsymbol\alpha_{\mathrm{cmd}},\qquad
\dot{\boldsymbol\omega}=\mathbf I^{-1}\mathbf B\mathbf u
\approx\boldsymbol\alpha_{\mathrm{cmd}} .
\label{eq:cascade_inner_dot}
\end{equation}
注意式中$\mathbf B\mathbf B^{-1}$只在同一局部工作点理想相消；上一拍工作点、转速来源不对称、
电机/舵机动态和限幅都会留下残差，故不能称“精确对角化”。

忽略积分项时，速率环输出角加速度
$\boldsymbol\alpha_{\mathrm{cmd}}=\mathbf K_{p,\omega}
(\boldsymbol\omega_{\mathrm{ref}}-\boldsymbol\omega)$，于是单通道理想传递函数为
\begin{equation}
G_i(s)=\frac{k_{p,\omega,i}}{s+k_{p,\omega,i}},
\label{eq:cascade_inner_tf}
\end{equation}
其名义带宽为$k_{p,\omega,i}$，而不是$k_\omega/I_i$；惯量已在力矩生成时乘入一次。
PI积分会增加一个状态，实际带宽与阻尼还受输出限幅和采样影响。

外环使用完整旋转向量$\mathbf e_R$。固定参考、小误差下
$\dot{\mathbf e}_R\approx-\boldsymbol\omega$，若内环足够快且
$\boldsymbol\omega_{\mathrm{ref}}=\mathbf K_{p,R}\mathbf e_R$，则
\begin{equation}
\dot{\mathbf e}_R\approx-\mathbf K_{p,R}\mathbf e_R,
\qquad \omega_{\mathrm{BW},o}\approx k_{p,R} .
\label{eq:cascade_outer_dot}
\end{equation}
这里没有原始四元数矢量的$1/2$因子。时间尺度分离应比较实测闭环频响，
不能由旧公式给出“94倍”或“364倍”结论；电机$\tau_m=0.28$\,s本身也表明理想静态分配
不能代表所有执行器动态。

\subsection{灵敏度函数与鲁棒性分析}
\label{sec:cascade_sensitivity}

\subsubsection{扰动灵敏度}

设内环存在交叉耦合扰动$\mathbf{d}$（如反扭矩投影），
角速度动力学为$\dot{\boldsymbol{\omega}}=\mathbf{I}^{-1}\mathbf{B}\mathbf{u}+\mathbf{d}$。
以内环参考$\boldsymbol{\omega}_{\mathrm{ref}}$为输入、扰动$\mathbf{d}$为扰动源，
角速度输出的扰动灵敏度函数为
\begin{equation}
    S_d(s) = \frac{\omega(s)}{d(s)} = \frac{1}{s+k_{p,\omega}}
    = \frac{\tau_i}{\tau_i s+1},\qquad \tau_i=\frac{1}{k_{p,\omega}}
    \label{eq:cascade_sens}
\end{equation}
其幅值$|S_d(j\omega)|$在低频段约为$1/k_{p,\omega}$——
内环增益$k_{p,\omega}$越大，对耦合扰动的抑制越强。
这从频域解释了为何提高内环带宽是耦合抑制的关键。

\subsubsection{对模型误差的鲁棒性}

在线$\mathbf{B}_{\mathrm{true}}$依赖模型参数（$k_T,k_Q,a,b$）。
设真实效能矩阵$\tilde{\mathbf{B}}=(\mathbf{I}+\boldsymbol{\Delta})\mathbf{B}_{\mathrm{true}}$，
$\boldsymbol{\Delta}$为乘性模型误差。则式(\ref{eq:cascade_inner_dot})变为
\begin{equation}
    \dot{\boldsymbol{\omega}}
    = \mathbf{I}^{-1}(\mathbf{I}+\boldsymbol{\Delta})
      \mathbf{K}_\omega(\boldsymbol{\omega}_{\mathrm{ref}}-\boldsymbol{\omega})
    \label{eq:cascade_model_err}
\end{equation}
交叉项残余为$\mathbf{I}^{-1}\boldsymbol{\Delta}\mathbf{K}_\omega
(\boldsymbol{\omega}_{\mathrm{ref}}-\boldsymbol{\omega})$。
$\mathbf{B}_{\mathrm{true}}$并非由推力、摆角和转速传感器在线测量得到。
固件使用指令摆角、开环转速模型估计以及MODEL-DEFAULT中的$k_T,k_Q,a,b$计算Jacobian，
故$\boldsymbol\Delta$同时包含参数误差、$\tau_m$失配、舵机跟踪误差和局部线性化残差。
即使单个参数误差为$\pm10\%$，也不能在未计算结构化误差传播和矩阵范数之前直接断言
$\|\boldsymbol\Delta\|=0.1$或得到相对SAS的百分比优势。

\subsection{奇异摄动与稳定性}
\label{sec:cascade_sp}

可用奇异摄动（singular perturbation）框架解释理想内外环的时间尺度分离。
引入小参数$\epsilon=\tau_i/\tau_o$，在忽略饱和、时变效能和执行器动态后写成：
\begin{align}
    \dot{\boldsymbol{\varepsilon}} &= \frac{1}{2}\boldsymbol{\omega}
    \label{eq:sp_slow}\\
    \epsilon\dot{\boldsymbol{\omega}} &= \frac{\tau_o}{I}\mathbf{K}_\omega
    \left(-\mathbf{K}_p\boldsymbol{\varepsilon}-\boldsymbol{\omega}\right)
    \label{eq:sp_fast}
\end{align}
其中$\boldsymbol{\varepsilon}$为慢变量、$\boldsymbol{\omega}$为快变量。
当$\epsilon\to0$时，快子系统在边界层内指数稳定于拟稳态
$\boldsymbol{\omega}=-\mathbf{K}_p\boldsymbol{\varepsilon}$，
慢子系统退化为$\dot{\boldsymbol{\varepsilon}}=-\frac{1}{2}\mathbf{K}_p\boldsymbol{\varepsilon}$。
只有进一步满足Tikhonov定理所要求的正则性、一致指数稳定性和初值条件时，
理想模型的解才可在相应时间区间内以$O(\epsilon)$逼近退化系统\cite{tikhonov1952}。
因此，减小$\epsilon$通常有利于理想模型的尺度分离，但不是现行固件精确解耦的保证。

需要强调的是，上述奇异摄动推导只为时间尺度分离提供近似解释。
当前多轴非线性固件还包含时变效能矩阵、饱和和执行器滞后，
因而不能由该推导直接得到完整闭环的$O(\epsilon)$误差界或稳定性结论。

\subsection{单通道Lyapunov检查与LMI鲁棒界}
\label{sec:cascade_lmi}

本节保留早期单通道LMI分析，并按当前固件变量重新审计。它只覆盖理想静态分配、
无执行器动态、无饱和的单通道线性近似；不构成多通道非线性固件的严格稳定性证明。
数值结果由\path{simulations/vector-thrust-lab/scripts/}目录中的三个Python脚本给出：
\begin{itemize}
  \item \texttt{cascade\_closedform.py}——符号推导与标称 $\mathbf{A}_0$ 自检；
  \item \texttt{robust\_lmi\_analysis.py}——LMI 求解（Petersen 引理、增益裕度）；
  \item \texttt{robust\_montecarlo.py}——蒙特卡洛数值验证。
\end{itemize}

\subsubsection{标称单通道闭环的Lyapunov检查}

单通道取状态$\mathbf{x}=[e_R,\omega]^T$（$e_R$为短弧旋转向量误差，
小角度下$e_R\approx\Delta\theta$，$\omega$为该通道角速度）。由式(\ref{eq:cascade_outer_dot})与(\ref{eq:cascade_inner_dot})，
标称闭环为线性系统
\begin{equation}
    \dot{\mathbf{x}} = \mathbf{A}_0\,\mathbf{x},\qquad
    \mathbf{A}_0 = \begin{bmatrix} 0 & -1 \\[2pt] k_{p,\omega}k_{p,R} & -k_{p,\omega} \end{bmatrix}
    \label{eq:cascade_A0}
\end{equation}
其特征多项式为
\begin{equation}
    \det(\lambda\mathbf{I}-\mathbf{A}_0)
    = \lambda^2+k_{p,\omega}\lambda+k_{p,\omega}k_{p,R}
    \label{eq:cascade_charpoly}
\end{equation}
当$k_{p,\omega}>0$且$k_{p,R}>0$时，该二阶多项式Hurwitz；是否过阻尼由
$k_{p,\omega}^2-4k_{p,\omega}k_{p,R}$决定。这一结果只属于上述理想单通道模型。

由 Lyapunov 第二方法，$\mathbf{A}_0$ Hurwitz 当且仅当存在$\mathbf{P}\succ 0$满足
\begin{equation}
    \mathbf{A}_0^T\mathbf{P}+\mathbf{P}\mathbf{A}_0 \prec \mathbf{0}
    \label{eq:cascade_lyap_lmi}
\end{equation}
此即线性矩阵不等式（LMI）。对仿真侧参数
$k_{p,R}=2.5$、$k_{p,\omega}=60$（源 \texttt{parameters.mjs}），
cvxpy/CLARABEL 按式(\ref{eq:cascade_A0})求得一组可行解
$\mathbf{P}=\bigl[\begin{smallmatrix}0.3193 & -0.0285 \\ -0.0285 & 0.0179\end{smallmatrix}\bigr]\succ 0$。
其特征值为$\{0.0153,\,0.3219\}$，
$\mathbf{A}_0^T\mathbf{P}+\mathbf{P}\mathbf{A}_0$的特征值为
$\{-10.5329,\,-0.1199\}$。由于$\mathbf P$并不唯一，以上数值只用于复现求解器证书；
可推广的结论仍限于式(\ref{eq:cascade_A0})所定义的理想单通道线性模型。

\subsubsection{参数摄动下的鲁棒界——Petersen 引理}

在线$\mathbf{B}_{\mathrm{true}}$依赖模型参数（$k_T,k_Q$），惯量$I$亦为组件级估算（\texttt{MODEL-DEFAULT}）。
设内环有效增益承受乘性摄动，集总为标量$\delta$：
$\tilde{\mathbf{A}} = \mathbf{A}_0 + \delta\,\mathbf{N}$，
其中$\mathbf{N}$作用于$\mathbf{A}_0$第二行（内环增益所在行）。写成范数有界形式
\begin{equation}
    \tilde{\mathbf{A}}(\delta) = \mathbf{A}_0 + \mathbf{B}_\Delta\,\delta\,\mathbf{C}_\Delta,\qquad
    \mathbf{B}_\Delta = \begin{bmatrix}0\\1\end{bmatrix},\quad
    \mathbf{C}_\Delta = \begin{bmatrix}\tfrac{K_\omega}{I}K_p, & -\tfrac{K_\omega}{I}\end{bmatrix}
    \label{eq:cascade_uncertainty}
\end{equation}
$\delta\approx\delta_B-\delta_I$（$\mathbf{B}$矩阵误差$-$惯量误差，一阶展开）。
利用Petersen引理~\cite{petersen1987stabilization}可构造二次稳定的充分LMI条件；
若存在$\mathbf P\succ0$、$\varepsilon>0$满足
\begin{equation}
    \begin{bmatrix}
      \mathbf{A}_0^T\mathbf{P}+\mathbf{P}\mathbf{A}_0+\varepsilon\,\mathbf{C}_\Delta^T\mathbf{C}_\Delta
        & \gamma\,\mathbf{P}\mathbf{B}_\Delta \\[3pt]
      \gamma\,\mathbf{B}_\Delta^T\mathbf{P} & -\varepsilon
    \end{bmatrix} \prec \mathbf{0}
    \label{eq:cascade_petersen}
\end{equation}
对固定$\gamma$这是 LMI；二分搜索$\gamma$得最大允许摄动界$\bar\gamma$。
它给出被证实稳定的保守摄动域，超过该域并不等于系统必然失稳。

\subsubsection{数值结果}

表\ref{tab:lmi_results} 给出按式(\ref{eq:cascade_A0})重新求解的
Lyapunov 证书与最大摄动界。结果可由
\path{simulations/vector-thrust-lab/scripts/robust_lmi_analysis.py}复现。

\begin{table}[htbp]
\centering
\caption{修正后单通道模型（式\ref{eq:cascade_A0}）的 LMI 结果}
\label{tab:lmi_results}
\small
\small
\begin{tabularx}{\textwidth}{>{\raggedright\arraybackslash}X >{\raggedright\arraybackslash}p{4.1cm}}
\toprule
\textbf{项} & \textbf{结果} \\
\midrule
\multicolumn{2}{l}{\textbf{仿真侧：}$k_{p,R}=2.5$，$k_{p,\omega}=60$} \\
Lyapunov 证书 & 存在 \\
最大允许摄动界 $\bar\gamma$ & $0.968$ \\
真实失稳边界 $\delta^\ast$ & $-1.000$（内环增益归零） \\
LMI 保守度 $\bar\gamma/|\delta^\ast|$ & $0.968$ \\
$|\delta|\le\bar\gamma$ 内的特征值抽样（5000点） & 未检出失稳样本 \\
\midrule
\multicolumn{2}{l}{\textbf{现行固件物理域等效：}$k_{p,R}=2.8$，$k_{p,\omega}=16.04$} \\
Lyapunov 证书；最大允许摄动界 & 存在；$\bar\gamma=0.897$ \\
\midrule
\multicolumn{2}{l}{\textbf{迁移失败反例（旧值未缩放）：}$k_{p,R}=2.8$，$k_{p,\omega}=0.28$} \\
Lyapunov 证书；最大允许摄动界 & 存在（欠阻尼）；$\bar\gamma=0.298$ \\
\bottomrule
\end{tabularx}
\end{table}

\begin{remark}
\textbf{固件速率环的现行量纲。}2026-08-11迁移后，PID输入为deg/s、输出为deg/s$^2$，
滚转/俯仰$k_{p,\omega}=16.042818$；进入惯量逆解前再乘$\pi/180$。
因此连续物理域的等效比例带宽仍为$16.04$\,rad/s，与迁移前
$0.28\times180/\pi$完全一致。表\ref{tab:lmi_results}中的$0.28$行不再代表当前固件，
而是“代码已增加$\pi/180$转换、持久化参数却仍保留旧值”的迁移失败反例；该故障会把
物理增益降低$57.3$倍，并把$\bar\gamma$由$0.897$变为$0.298$。
现行参数下时间尺度分离比为$16.04/2.8=5.7$，
达到常用的$3$--$5$倍经验量级，且$\bar\gamma=0.897$。
真实失稳机制为$\delta\to-1$时内环有效增益$\to0$（系统退化为开环临界），
Petersen 引理在单标量不确定上近紧（保守度$3.2\%$（仿真）$\sim10.3\%$（固件））。
该界只覆盖理想静态分配、无执行器动态、无饱和的单通道线性近似，
不能覆盖多轴耦合、时变工作点和执行器滞后。
\end{remark}

\subsubsection{蒙特卡洛数值验证}

历史脚本对$\delta\in[-2,2]$做随机采样（每组$N=3000\sim5000$）：
$\bar\gamma$内5000个旧模型样本未出现失稳；
$\bar\gamma$外（$\delta\in[-1.017,-0.968]$，越过$\delta^\ast=-1$）失稳率$35.3\%$。
特征值判定与时域积分（$8\,\mathrm{s}$，精确矩阵指数）在所有测试点一致，
确认线性模型自洽。详见 \texttt{robust\_montecarlo.py}。

\begin{remark}
\textbf{本节结论的边界与适用范围。}
(i) 上述证明在"$\mathbf{B}_{\mathrm{true}}\mathbf{B}_{\mathrm{true}}^{-1}=\mathbf{I}$精确对角化"假设下成立，
即三通道独立分析；多通道交叉耦合（$\boldsymbol{\Delta}$的非对角部分）需要重新构造
复合 Lyapunov 函数并显式界定通道间交叉项，不能由本节的三个标量证书直接推出。
(ii) $\bar\gamma\approx0.97$只描述旧标量模型，不能据此判定现行架构的主导风险，
也不能解释两次实机$k_Q$调参结果；
(iii) 但本模型\emph{不覆盖}第\ref{sec:cascade_workpoint}节"瞬态失配"类不确定——
即差速分配器用指令油门、力矩用实际转速（$\tau_m$ 滞后）导致的\textbf{时变}工作点失配。
实机 yaw 在$\tau_m$观测器接入前后的稳定差异（第\ref{sec:cascade_implementation}节）
即属此类，需要执行器动态增广的不确定模型（$\tau_m$作为未建模动态），超出单通道静态 LMI 范围。
(iv) $\mathbf{A}_0$ 使用 \texttt{MODEL-DEFAULT} 标称$I$值，台架标定后需重解。
\end{remark}

\subsection{相位裕度与差速通道的敏感性}
\label{sec:cascade_pm}

本节保留旧单通道频域模型，用于检查电机滞后和陀螺滤波可能消耗的相位裕度（PM）。
其增益采用仿真侧$K_\omega/I=60$的量纲；当前固件滚转/俯仰比例项在物理域的
等效带宽为$16.04$\,rad/s，两者不是同一整定点。
项目记录中的滤波相位阈值约为$4.2^\circ$和$6.9^\circ$；下述旧模型数值不能解释该实机现象。
数值由\path{phase_margin_analysis.py}给出。

\subsubsection{双通道开环传函}

级联架构的两类执行器时间常数截然不同，须分通道建立开环传函。
单通道内环（控制器测量滤波后角速度$\omega_f$，$\omega_{\mathrm{des}}=K_p e_R$）：

\textbf{摆角通道}（$\delta_f,\delta_t$，舵机响应快，执行器动态可略）：
\begin{equation}
    L_\delta(s) = \frac{K_\omega/I}{s\,(\tau_f s + 1)}
    \label{eq:pm_servo}
\end{equation}

\textbf{差速通道}（$\Delta\omega$，经电机一阶滞后$\tau_m$）。
将控制分配$\Delta\omega_{\mathrm{cmd}}=M_{\mathrm{des}}/B_{\mathrm{yaw}}$、
电机跟踪$\Delta\omega_{\mathrm{act}}=\Delta\omega_{\mathrm{cmd}}/(\tau_m s+1)$、
力矩$M=B_{\mathrm{yaw}}\Delta\omega_{\mathrm{act}}$、$\dot\omega=M/I_x$
四式联立，$B_{\mathrm{yaw}}/I_x$ 完全约去：
\begin{equation}
    L_\omega(s) = \frac{K_\omega/I}{s\,(\tau_f s + 1)(\tau_m s + 1)}
    \label{eq:pm_diff}
\end{equation}
其中$\tau_f$ 为陀螺一阶滤波的连续等效时间常数
（$\alpha=0.4$、$dt=0.005$\,s 下$\tau_f=-dt/\ln(1-\alpha)\approx 9.8$\,ms），
$\tau_m=0.28$\,s 为电机时间常数。

\subsubsection{历史数值结果：差速通道相位更敏感}

表\ref{tab:pm_compare} 给出两通道的增益穿越频率$\omega_{\mathrm{gc}}$与相位裕度。
在该旧模型中，差速通道的额外PM损失主要来自执行器$\tau_m$。

\begin{table}[htbp]
\centering
\caption{内环开环相位裕度对比（仿真侧 \texttt{parameters.mjs} 增益 $K_\omega/I=60$；固件等效带宽需按当前PI另行复核）}
\label{tab:pm_compare}
\begin{tabular}{lcc}
\hline
 & 摆角通道 $L_\delta$ & 差速通道 $L_\omega$ \\
\hline
增益穿越频率 $\omega_{\mathrm{gc}}$ & $53.2$\,rad/s ($8.5$\,Hz) & $14.4$\,rad/s ($2.3$\,Hz) \\
相位 @$\omega_{\mathrm{gc}}$：积分器 $1/s$ & $-90^\circ$ & $-90^\circ$ \\
\quad 执行器 $1/(\tau_m s+1)$ & --- & $-76^\circ$ \\
\quad 滤波 $1/(\tau_f s+1)$ & $-27.5^\circ$ & $-8.0^\circ$ \\
\quad 合计 & $-117.5^\circ$ & $-174.0^\circ$ \\
\textbf{相位裕度 PM} & $\mathbf{62.5^\circ}$（充足） & $\mathbf{6.0^\circ}$（临界） \\
增益裕度 GM & $\infty$ & $1.76$ \\
\hline
\end{tabular}
\end{table}

差速内环的闭环特征方程
$\tau_f\tau_m s^3+(\tau_f+\tau_m)s^2+s+K_\omega/I=0$，
其主导共轭极点为$-0.74\pm 14.5j$，阻尼比$\zeta=0.051$——
与频域 PM 的经验关系$\mathrm{PM}\approx 100\zeta$ 一致（$100\zeta=5.1^\circ \approx \mathrm{PM}=6.0^\circ$），
构成频域–时域双重交叉验证。

\subsubsection{滤波$\alpha$ 扫描与失稳边界}

表\ref{tab:pm_alpha} 扫描陀螺滤波强度$\alpha$（$\alpha$ 越小、滤波越重、$\tau_f$ 越大）。
旧模型中摆角通道在扫描区间保持正PM，差速通道在$\alpha\le0.2$时转为负PM。
该边界按固件量纲的偏移见下文"固件量纲复核"子节（机制不变，临界$\alpha$数值随带宽漂移）。

\begin{table}[htbp]
\centering
\caption{滤波$\alpha$ 扫描下的相位裕度（仿真侧 \texttt{parameters.mjs} 现行增益）}
\label{tab:pm_alpha}
\begin{tabular}{ccccc}
\hline
$\alpha$ & $\tau_f$ (ms) & 摆角 PM & 差速 PM & 判定 \\
\hline
$0.99$ & $1.1$ & $86.3^\circ$ & $13.0^\circ$ & 稳 \\
$0.80$ & $3.1$ & $79.6^\circ$ & $11.4^\circ$ & 稳 \\
$0.40$（标称） & $9.8$ & $62.5^\circ$ & $6.0^\circ$ & 临界 \\
$0.20$ & $22.4$ & $46.0^\circ$ & $-3.3^\circ$ & 失稳 \\
$0.10$ & $47.5$ & $32.9^\circ$ & $-17.0^\circ$ & 失稳 \\
\hline
\end{tabular}
\end{table}

\begin{remark}
\textbf{参数来源与结论边界。}本节的相位裕度数值基于\textbf{仿真侧} \texttt{parameters.mjs}
现行增益（$K_\omega/I=60$）。固件内环按速率环增益 $k_{p,\omega}$ 直接定义名义带宽
（式\ref{eq:cascade_inner_tf}，惯量已在力矩生成时乘入）；量纲迁移后物理域等效带宽
为$16.04$\,rad/s。可以保留的机制性判断是：相较摆座通道，差速通道多一个
电机动态环节；同一穿越频率下，该环节会增加相位滞后，陀螺滤波还会继续消耗裕度。
但穿越频率本身随控制增益和对象增益变化，故不能说PM只由$\tau_m$决定。

旧模型的PM约$6^\circ$与项目抖振记录数值接近，只能作为待检验假设。
$\tau_m$开环模型估计器改善的是B矩阵工作点匹配，并未消除电机的物理滞后，
不能直接等价为PM提升。要建立实机解释，需从当前PI、采样保持、滤波器、BTRUE调度和
电机动态推导开环模型，再用扫频或闭环辨识验证。
\end{remark}

\subsubsection{固件量纲复核：PM 对带宽的依赖}

上述质疑 (i) 成立：PM 的绝对值由增益穿越频率决定，而穿越频率随控制增益
$k_{p,\omega}$ 与对象增益变化，不能宣称固定 PM。表\ref{tab:pm_K_dep} 给出
差速通道 PM 对 $k_{p,\omega}$ 的依赖（$\tau_m=0.28$\,s、$\tau_f=9.8$\,ms，
\texttt{scripts/pm\_sensitivity.py}）：
\begin{table}[htbp]
\centering
\caption{差速通道 PM 对等效内环带宽的依赖（$\tau_m=0.28$\,s、$\tau_f=9.8$\,ms）}
\label{tab:pm_K_dep}
\begin{tabular}{ccc}
\hline
带宽来源 & 等效 $k_{p,\omega}$ (rad/s) & 差速 PM \\
\hline
迁移失败反例（旧参数未缩放） & $0.28$ & $85.4^\circ$（物理增益降低57.3倍） \\
现行固件物理域等效 & $16.04$ & $22.5^\circ$（偏紧） \\
仿真 \texttt{parameters.mjs} & $60$ & $6.0^\circ$（临界） \\
旧论文假设（已废止） & $364$ & $-13.1^\circ$（必失稳） \\
\hline
\end{tabular}
\end{table}

几点说明：
\begin{enumerate}
  \item \textbf{量纲迁移后的等效性}：现行固件以$16.042818$接收deg/s误差并输出
        deg/s$^2$，随后乘$\pi/180$进入物理层，等效带宽为$16.04$\,rad/s。
        表中的$0.28$只模拟旧参数被错误保留的升级故障；它不是另一种合理整定，也不是
        当前代码的“名义值”。第\ref{sec:cascade_lmi}节的$0.897$与$0.298$对照说明
        参数迁移和持久化版本控制会直接改变稳定性分析结论。
  \item \textbf{质疑 (ii) 成立}：$\tau_m$ 开环观测器消除的是\textbf{瞬态工作点失配}
        （B 矩阵用 $w_{\mathrm{est}}\approx w_{\mathrm{act}}$ 而非指令油门，
        有效增益不再随油门瞬态漂移），\emph{不是}电机的物理滞后。因此它恢复的是
        "PM 不瞬态崩塌"，而非"PM 数值提升"；物理滞后仍由第\ref{sec:cascade_workpoint}节的
        实际回路承担。数值上，无观测器时油门释放瞬间有效增益 $\times(w_{\mathrm{act}}/w_{\mathrm{cmd}})^2$
        在$k_{p,\omega}=60$的仿真侧示例中可达 $\sim2.8$ 倍（等效
        $k_{p,\omega}\to166$，简化模型 PM 瞬态 $\to-4.7^\circ$，负裕度窗口约 200\,ms）。
        该结果用于说明工作点失配机制；观测器只能减小模型工作点误差，不能保证真实有效增益恒定。
  \item \textbf{机制性结论稳健}：无论采用哪组带宽，差速通道因多一个 $\tau_m$ 动态
        环节而相位裕度低于摆角通道、且随滤波 $\alpha$ 收紧单调下降。临界点依赖带宽：
        表\ref{tab:pm_alpha}的$k_{p,\omega}=60$模型在$\alpha\le0.2$时已为负裕度，
        而$k_{p,\omega}=16.04$的等效比例模型在$\alpha=0.2$时仍约为$17.6^\circ$，
        直至$\alpha\approx0.05$才转为负值。
        绝对值须由扫频/闭环辨识实测校准——本节的贡献是给出可复算的依赖关系与量纲换算，
        而非未经验证的 PM 数值。
\end{enumerate}

\subsection{控制参数优化与离线整定}
\label{sec:cascade_tuning}

级联架构有三个待整定增益（$K_p,K_i,K_\omega$）。
本节记录历史离线网格搜索方法；“最优”仅指所选离散网格和复合代价函数的最小值。

\subsubsection{多目标代价函数}

参数整定的目标是在耦合抑制、跟踪速度与执行器代价间取得平衡。
定义复合代价函数
\begin{equation}
    J(\mathbf{g}) = w_1\,\theta_{\mathrm{peak}} + w_2\,t_{\mathrm{settle}}
    + w_3\,\bar{u}
    \label{eq:cascade_cost}
\end{equation}
其中$\mathbf{g}=[K_p,K_i,K_\omega]^T$，
$\theta_{\mathrm{peak}}$为耦合扰动下的俯仰峰值，
$t_{\mathrm{settle}}$为调节时间（回到峰值$10\%$内），
$\bar{u}$为控制能量（执行器指令幅值均值），
$w_1,w_2,w_3$为权重。

\subsubsection{网格搜索与精化}

采用两阶段离线搜索：
\begin{enumerate}[label=(\arabic*)]
    \item \textbf{粗网格}：在$K_p\in[1,3]$、$K_i\in[0,1]$、$K_\omega\in[4,12]$
          上逐点仿真评估$J$；
    \item \textbf{精化}：在最优点附近细分网格，重复搜索至收敛。
\end{enumerate}
历史脚本中的$J$曲面在低增益区较平缓（图\ref{fig:cascade_param_opt}），
网格最小点为$K_p\approx1.0$、$K_i\approx0$、$K_\omega\approx4.0$。
该点未按当前固件PI量纲、$15^\circ$限幅和电机观测器重算，不能作为实机整定值。

\begin{figure}[H]
\centering
\includegraphics[width=0.9\textwidth]{fig-sim/cascade_param_opt.pdf}
\caption{历史级联参数网格的代价曲面。红星为该离散网格与旧代价定义下的最小点。}
\label{fig:cascade_param_opt}
\end{figure}

\subsubsection{整定原则的物理理解}

历史搜索提示，在线效能矩阵可在一定程度上降低对高反馈增益的依赖，
但交叉项残差仍受Jacobian失配和饱和影响。提高$K_\omega$通常加快内环并增强低频扰动抑制，
同时也会提高噪声与未建模动态敏感性；过低增益则削弱时间尺度分离。
因此实机整定仍需在跟踪、噪声、相位裕度和执行器代价之间折中。
$K_i$的引入主要用于消除常值扰动的稳态误差，但会略微增大耦合峰值
（积分项的慢放电与耦合瞬态叠加），故在耦合抑制优先的场景可取$K_i=0$。

\subsection{仿真验证与可视化}
\label{sec:cascade_sim}

本节图表为2026年7月历史资产，包含$25^\circ$场景，超过当前固件$\pm15^\circ$限幅；
LQR、INDI与级联脚本还使用不同控制器定义和工作点。为保留迭代过程，原始数据继续列出，
但不再用于“精确解耦”“两个数量级改善”或控制器排名，需由统一当前管线重新生成。

\subsubsection{交叉耦合时域响应}

在历史局部配平模型中，施加上摆$\delta_f=25^\circ$阶跃，
图\ref{fig:cascade_time}对比了LQR、INDI与级联架构的三通道时域响应。

\begin{figure}[H]
\centering
\includegraphics[width=0.82\textwidth]{fig-sim/cascade_time_response.pdf}
\caption{历史脚本的$\delta_f=25^\circ$阶跃响应。该摆角超过固件限幅，
三种控制器的数据口径也未统一。}
\label{fig:cascade_time}
\end{figure}

旧脚本记录的级联、LQR和INDI峰值见表\ref{tab:cascade_ch_coupling2}。
由于同类LQR基线在相邻章节出现$1.33^\circ$与$12.05^\circ$的不一致，
本表只能作为数据管线审计材料，不能据数值差异解释控制机理。

\begin{table}[H]
\centering
\caption{历史脚本的俯仰耦合峰值（单位：度，5次运行均值）。
该表与表\ref{tab:lqr_coupling_deg}的基线不一致，仅作为数据管线审计记录。}
\label{tab:cascade_ch_coupling2}
\small
\begin{tabular}{l c c}
\toprule
\textbf{控制架构} & $\delta_f=10^\circ$ & $\delta_f=25^\circ$ \\
\midrule
LQR（冻结$\mathbf{B}$，单级） & 12.05 & 5.91 \\
INDI（在线$\mathbf{B}$，单级） & 35.86 & 41.51 \\
\textbf{级联P外环 + 在线$\mathbf{B}_{\mathrm{true}}$内环} & \textbf{0.06} & \textbf{0.08} \\
\bottomrule
\end{tabular}
\end{table}

\subsubsection{噪声鲁棒性}

扫描陀螺噪声标准差，对比三架构的耦合峰值（图\ref{fig:cascade_noise}）。

\begin{figure}[H]
\centering
\includegraphics[width=0.72\textwidth]{fig-sim/cascade_noise_robust.pdf}
\caption{历史脚本的噪声扫描。反常趋势与跨表基线不一致使其不能用于鲁棒性排序。}
\label{fig:cascade_noise}
\end{figure}

\begin{table}[H]
\centering
\caption{历史脚本在不同陀螺噪声下的俯仰耦合峰值（单位：度，5次运行均值）}
\label{tab:cascade_noise}
\small
\begin{tabular}{l c c c c c}
\toprule
\textbf{架构} & 0 & 0.005 & 0.01 & 0.02 & 0.04 rad/s \\
\midrule
LQR & 12.05 & 12.04 & 12.04 & 12.04 & 12.03 \\
INDI & 36.00 & 36.03 & 37.28 & 33.41 & 28.80 \\
\textbf{级联} & \textbf{0.05} & \textbf{0.10} & \textbf{0.16} & \textbf{0.24} & \textbf{0.42} \\
\bottomrule
\end{tabular}
\end{table}

该历史表内部也存在反常趋势：INDI所谓“噪声退化”从$36.0^\circ$下降到$28.8^\circ$，
LQR基线在不同表中从$1.33^\circ$变为$12.05^\circ$。这些差异不能用“口径不同”替代统一实验，
因此本表只作为数据管线不一致的审计证据，不支持三架构噪声性能排序。

\subsubsection{执行器速率饱和}

扫描内环增量限幅$\Delta u_{\max}$，评估执行器速率饱和对耦合抑制的影响
（图\ref{fig:cascade_sat}）。

\begin{figure}[H]
\centering
\includegraphics[width=0.68\textwidth]{fig-sim/cascade_saturation.pdf}
\caption{历史增量原型中的执行器速率限幅扫描。该图用于记录旧实验口径，
不代表当前绝对指令固件的饱和鲁棒性。}
\label{fig:cascade_sat}
\end{figure}

图\ref{fig:cascade_sat}来自历史增量原型，而固件采用绝对指令并实施幅值限位；
因此不能把$\Delta u_{\max}$扫描解释为当前执行器速率饱和鲁棒性。

\subsubsection{Tail-sitter大姿态角验证}

在tail-sitter垂直起飞（3秒slerp过渡至$\theta=-90^\circ$）下，
四元数表示在万向锁附近保持正则；表中的精确误差值属于历史脚本，
控制器与当前固件未对齐，不能用于优劣排序或性能验证。

\begin{table}[H]
\centering
\caption{历史tail-sitter脚本的3\,s slerp姿态误差输出}
\label{tab:cascade_ch_vtol2}
\small
\begin{tabular}{l c c}
\toprule
\textbf{控制架构} & $\|\boldsymbol{\varepsilon}_{\mathrm{err}}\|$峰值 & 末值（10\,s） \\
\midrule
SAS（增益调度+前馈） & 0.0265 & $4.9\times10^{-4}$ \\
INDI（v3，含惯量逆） & 0.0355 & $9.1\times10^{-4}$ \\
LQR & 0.0087 & $9.5\times10^{-6}$ \\
\textbf{级联（误差四元数外环）} & 0.0087 & $1.4\times10^{-3}$ \\
\bottomrule
\end{tabular}
\end{table}

\subsection{分析讨论}
\label{sec:cascade_analysis}

\subsubsection{三种"LQR在线B"路径的失败教训}

曾系统评估三种LQR交叉项补偿在线化的直接路径
（增益调度LQR、SDRE在线Riccati、LQR外环+在线B分配）。
旧脚本报告三条路径的峰值为$1.3^\circ$--$2.6^\circ$，但没有统一实现、统计区间和
消融试验，无法把差异唯一归因于“单级全状态反馈”。这些负结果只说明直接替换效能矩阵
没有在该原型中自动改善指标；控制架构、权重、饱和和工作点的因果贡献仍需分别验证。

\subsubsection{架构的适用边界}

\begin{enumerate}[label=(\arabic*)]
    \item \textbf{时间尺度分离}：内环带宽须足够高。执行器响应迟缓时解耦退化；
    \item \textbf{低油门退化}：$\omega_0\to0$时$\mathbf{B}_{\mathrm{true}}$奇异，
          需回退对角分配或舵面控制；
    \item \textbf{执行器饱和}：大幅扰动下摆角和差速限幅会破坏线性分配关系；
          当前证据不足以在绝对分配与INDI增量分配之间给出统一的饱和性能排序。
\end{enumerate}

\subsubsection{与其他架构的定位关系}

\begin{itemize}
    \item LQR适合在明确工作点与二次型指标下进行系统化线性反馈设计；
    \item INDI适合研究具备高质量角加速度估计时的增量动态补偿；
    \item 需固定计算量、在线局部增益与易调试数据流：当前级联架构具有工程适配性；
          在统一对比完成前不作“首选”或性能排名。
\end{itemize}

\subsection{结论}
\label{sec:cascade_conclusion}

本章统一了固件五层链路：四元数短弧外环、速率PI、惯量与陀螺前馈、任务轴置换以及
上一拍Jacobian的绝对准线性分配。修正后的单通道分析说明该结构可在理想局部条件下形成
带宽分离，但不构成多通道非线性系统的精确对角化或严格稳定性证明。
历史图表因控制器定义、工作点、限幅和数据管线不一致，只作为审计材料保留；
“两个数量级改善”、全噪声区间最优和饱和鲁棒性不再作为结论。
后续需按当前固件数据流重建统一仿真，将PI积分、离散采样、饱和和执行器动态纳入鲁棒分析，
并完成HIL与带原始日志的重复飞行试验。

\subsection{实机实现状态与工程教训（2026-08）}
\label{sec:cascade_implementation}

本章架构已于 2026-08 在实机飞控（STM32H743）上线，本节记录实现状态与
调试中沉淀的工程记录。这些记录支持链路部署与问题定位，但不能替代闭环理论证明或性能试验。

\subsubsection{实现状态}

\begin{itemize}
    \item \textbf{五层链路全线上}：参考姿态生成（三模式）→ 四元数几何外环
          （精确角度提取 + 短弧）→ 速率 PI → 惯量逆解 + 前馈 → 在线 B 分配；
    \item \textbf{垂直/水平位置环已接入}：定高（位置 P + 速度 PI + 油门杆=高度率指令 +
          GUIDED 加速度直通）、定点（Loiter 状态机 + 物理刹车预测 + 位置/速度串级），
          与姿态五层链路级联；
    \item \textbf{FPV 摇杆曲线}（RATE\_MODE）：三参数模型实装，默认满杆
          roll/pitch 333°/s、yaw 222°/s（第\ref{sec:rc_curve}节）；
    \item \textbf{在线辨识}：RLS为纯观测模式；项目记录含$b$、$d$与$k_Q/k_T$调参值，
          但论文工程缺少原始日志、单位和重复次数，故不作为物理标定结论（第\ref{sec:online_id}节）。
\end{itemize}

\subsubsection{工程教训（理论未覆盖的实践修正）}

\begin{enumerate}[label=(\arabic*)]
    \item \textbf{限幅必须解耦}：RATE\_MODE 限幅最初复用姿态外环 80°/s
          （$k_{\mathrm{MaxTargetRate}}$）——曲线在约半杆即被削平，满杆暴力手感
          完全消失。2026-08-11 独立为 600°/s（曲线理论值 333°/s 之上留余量），
          限幅从"姿态环约束"变为"曲线之上的安全网"（第\ref{sec:rc_curve}节）；
    \item \textbf{工作点观测是分配器的关键条件}：$\tau_m$模型观测器未接入前，油门瞬态下
          B 矩阵工作点与实际转速失配，有效增益瞬态越震荡点约 160\,ms
          激发yaw振铃；模型预测器接入后项目记录显示振铃减弱，但不能宣称瞬态增益恒定
          （2026-08-09 三层定位，第\ref{sec:allocation}章）；
    \item \textbf{通道积分独立性}：2026-08-10 通读发现 pitch 内环误用 roll 的
          PID 实例（复制粘贴 bug）——pitch 积分恒 0 且双通道积分互相污染，
          实机表现为"pitch 一直飘、积分无法消除"；修复后积分独立。
          教训：同构三通道代码逐一核对实例名，勿因参数相同想当然；
    \item \textbf{数值病态区必须显式规避}：B 各项 $\propto w_0^2$、det $\propto w_0^6$，
          低油门工作点冻结（$0.6\,w_{\mathrm{hover}}$）+ 零油门门控是实机可飞的前提，
          纯理论推导不会暴露此问题。
\end{enumerate}

\subsubsection{与 NDI/INDI 的定位关系（2026-08-11 精确归类）}

本架构为\textbf{NDI型绝对准线性分配}（绝对指令、模型前馈、在线局部B和工作点估计），
\textbf{非 INDI}（无增量累加、无实测角加速度）。三条实现路径（SAS/INDI/LQR）
的教训与本章架构的定位：
\begin{itemize}
    \item SAS：对角分配忽略交叉项——本架构以在线局部B减小该近似误差；
    \item INDI：增量形式对模型误差鲁棒，但依赖实测角加速度（噪声放大）——
          本架构以$\tau_m$模型观测器与反馈积分降低部分失配，但尚不能据此宣称实现同等效果；
    \item LQR：冻结B对工作点变化敏感——本架构用在线局部B与模型工作点估计缓解，而非消除漂移。
\end{itemize}
